
\documentclass[9pt, headings=standardclasses, parskip=half]{scrartcl}
\usepackage{ifthen}
\usepackage{iftex}
\usepackage{csquotes}

\usepackage[T1]{fontenc}     % handle accented glyphs properly
\usepackage[english]{babel}  % load the hyphenation patterns you need

\usepackage[automark]{scrlayer-scrpage}
% \clearpairofpagestyles
% \ofoot{\pagemark} % für ein einseitiges Dokument: \ofoot platziert den Inhalt in der äußeren Fußzeile (unten rechts)
\pagestyle{scrheadings}

\usepackage{graphicx}
\usepackage[dvipsnames]{xcolor}
\usepackage{soulutf8}
\usepackage{xparse}

\ExplSyntaxOn
\tl_new:N \__l_SOUL_argument_tl
\cs_set_eq:Nc \SOUL_start:n { SOUL@start }
\cs_generate_variant:Nn \SOUL_start:n { V }
\cs_set_protected:cpn {SOUL@start} #1
 {
  \tl_set:Nn \__l_SOUL_argument_tl { #1 }
  \regex_replace_all:nnN
   { \c{\(} (.*?) \c{\)} } % look for \(...\) (lazily)
   { \cM\$ \1 \cM\$ }      % replace with $...$
   \__l_SOUL_argument_tl
  \SOUL_start:V \__l_SOUL_argument_tl % do the usual
 }
\ExplSyntaxOff

\let\SOULhl\hl % save soul's \hl under a private name
\renewcommand{\hl}[1]{\SOULhl{#1}} % keep plain \hl working if you want

\definecolor{HLgreen}{HTML}{77DD77}
\definecolor{HLyellow}{HTML}{FFFF66}
\definecolor{HLorange}{HTML}{FFB347}
\definecolor{HLred}{HTML}{FF6961}
\definecolor{HLpink}{HTML}{FFB6C1}
\definecolor{HLturquoise}{HTML}{40E0D0}

\RenewDocumentCommand{\hl}{O{1} m}{% % master highlight macro: \hl[level]{text}, default = green
  \begingroup
  \IfEqCase{#1}{%
    {1}{\sethlcolor{HLgreen}\SOULhl{#2}}%
    {2}{\sethlcolor{HLyellow}\SOULhl{#2}}%
    {3}{\sethlcolor{HLorange}\SOULhl{#2}}%
    {4}{\sethlcolor{HLred}\SOULhl{#2}}%
    {5}{\sethlcolor{red}\SOULhl{#2}}%
    {6}{\sethlcolor{HLturquoise}\SOULhl{#2}}%
  }[\PackageError{hl}{Undefined highlight level: #1}{}]%
  \endgroup
}

\usepackage[left=38mm, right=38mm, top=20mm, bottom=30mm]{geometry}

\usepackage{amsmath, amssymb, amsthm}
\usepackage{mathtools}
\usepackage{mathdots} % without this package, only \vdots and \ddots are taken from the text font (not the math font), which looks bad if the text font is different from the math font
\usepackage{accents}

\usepackage{enumitem}
\renewcommand{\labelitemi}{\textbullet}
\renewcommand{\labelitemii}{\raisebox{0.1ex}{\scalebox{0.8}{\textbullet}}}
\renewcommand{\labelitemiii}{\raisebox{0.2ex}{\scalebox{0.6}{\textbullet}}}
\renewcommand{\labelitemiv}{\raisebox{0.3ex}{\scalebox{0.4}{\textbullet}}}

\usepackage{caption, subcaption}

\usepackage[backend=biber,style=numeric,sorting=none]{biblatex}
\addbibresource{\jobname.bib}

\usepackage{pifont}

\usepackage{tikz}
\usetikzlibrary{arrows, arrows.meta, shapes, positioning, calc, fit, patterns, intersections, math, 3d, tikzmark, decorations.pathreplacing, decorations.markings}
\usepackage{tikz-dependency, tikz-qtree, tikz-qtree-compat}
\usepackage{tikz-3dplot}
\usepackage{tikzpagenodes}
% \usetikzlibrary{external}
% \tikzexternalize[prefix=tikz/]
\usepackage{pgfplots}

% define colors
\definecolor{funblue}{rgb}{0.10, 0.35, 0.66}
\definecolor{alizarincrimsonred}{rgb}{0.85, 0.17, 0.11}
\definecolor{amethyst}{rgb}{0.6, 0.4, 0.8}
\definecolor{mypurple}{rgb}{128, 0, 128} 
\definecolor{highlightpurple}{rgb}{0.6, 0.0, 0.6}

\renewcommand{\emph}[1]{\textcolor{highlightpurple}{\textsl{#1}}}

\usepackage{thmtools}
\newlength{\thmspace}
\setlength{\thmspace}{3pt plus 1pt minus 1pt}
\declaretheoremstyle[headfont=\bfseries, bodyfont=\normalfont, spaceabove=\thmspace, spacebelow=\thmspace, qed=\ensuremath{\vartriangleleft}, postheadspace=1em]{assertionstyle}
\declaretheorem[style=assertionstyle, name=Theorem, numberwithin=section]{theorem} % <-- resets in each section
\declaretheorem[style=assertionstyle, name=Lemma,       sibling=theorem]{lemma}
\declaretheorem[style=assertionstyle, name=Corollary,   sibling=theorem]{corollary}
\declaretheorem[style=assertionstyle, name=Proposition, sibling=theorem]{proposition}
\declaretheorem[style=assertionstyle, name=Conjecture,  sibling=theorem]{conjecture}
\declaretheorem[style=assertionstyle, name=Claim,       sibling=theorem]{claim}
\declaretheorem[style=assertionstyle, name=Fact,        sibling=theorem]{fact}
\declaretheoremstyle[headfont=\bfseries, bodyfont=\normalfont, spaceabove=\thmspace, spacebelow=\thmspace, qed=\ensuremath{\blacktriangleleft}, postheadspace=1em]{definitionstyle}
\declaretheorem[style=definitionstyle, name=Definition, numberwithin=section]{definition}
\declaretheorem[style=definitionstyle, name=Problem,    sibling=definition]{problem}
\declaretheoremstyle[headfont=\bfseries, bodyfont=\normalfont, spaceabove=\thmspace, spacebelow=\thmspace, qed=\ding{45}, postheadspace=1em]{exercisestyle}
\declaretheorem[style=exercisestyle, name=Exercise, numberwithin=section]{exercise}
\declaretheoremstyle[headfont=\bfseries\color{red}, bodyfont=\normalfont, spaceabove=\thmspace, spacebelow=\thmspace, qed=\ensuremath{\color{red}\blacktriangleleft}, postheadspace=1em]{solutionstyle}
\declaretheorem[style=solutionstyle, name=Solution, numbered=no]{solution}
\declaretheoremstyle[headfont=\bfseries, bodyfont=\normalfont, spaceabove=6pt, spacebelow=6pt, qed=\ensuremath{\square}, postheadspace=1em]{proofstyle}
\let\proof\relax
\let\endproof\relax
\declaretheorem[style=proofstyle, name=Proof,    numbered=no]{proof}
\declaretheoremstyle[headfont=\bfseries, bodyfont=\normalfont\normalsize, spaceabove=\thmspace, spacebelow=\thmspace, qed=\ensuremath{\blacktriangleleft}, postheadspace=1em]{remarkstyle}
\declaretheorem[style=remarkstyle, name=Remark,   numberwithin=section]{remark}
\declaretheoremstyle[headfont=\bfseries\color{funblue}, bodyfont=\normalfont\normalsize, spaceabove=\thmspace, spacebelow=\thmspace, qed=\ensuremath{\color{funblue}\blacktriangleleft}, postheadspace=1em]{examplestyle}
\declaretheorem[style=examplestyle, name=Example, sibling=remark]{example}
\declaretheoremstyle[headfont=\color{alizarincrimsonred}\bfseries, bodyfont=\normalfont\normalsize, spaceabove=\thmspace, spacebelow=\thmspace, qed=\ensuremath{\color{alizarincrimsonred}\blacktriangleleft}, postheadspace=1em]{cautionstyle}
\declaretheorem[style=cautionstyle, name=Caution, sibling=remark]{caution}
\declaretheoremstyle[headfont=\bfseries, bodyfont=\normalfont\footnotesize, spaceabove=\thmspace, spacebelow=\thmspace, postheadspace=1em]{smallremarkstyle}
\declaretheorem[style=smallremarkstyle, name=Remark, sibling=remark]{smallremark}
\declaretheoremstyle[headfont=\bfseries\color{amethyst}, bodyfont=\normalfont, spaceabove=6pt, spacebelow=6pt, qed=\ensuremath{\color{amethyst}\blacktriangleleft}, postheadspace=1em]{digressionstyle}
\declaretheorem[style=digressionstyle, name=Digression, sibling=remark]{digression}

\numberwithin{equation}{section} % equations numbered within sections

\newenvironment{verticalhack}{\begin{array}[b]{@{}c@{}}\displaystyle}{\\\noalign{\hrule height0pt}\end{array}} % to make qed symbol aligned

\usepackage{algorithm}
\usepackage{algorithmicx}
\usepackage[italicComments=false]{algpseudocodex} % macht probleme mit TikZ externalization

\newcommand{\im}{\operatorname{Im}}

\DeclareMathOperator*{\argmax}{arg\,max}
\DeclareMathOperator*{\argmin}{arg\,min}
\DeclareMathOperator{\Span}{span}
\DeclareMathOperator{\Kern}{kern}
\DeclareMathOperator{\Trace}{trace}
\DeclareMathOperator{\Rank}{rank}

\newcommand*\matrspace{0.8mu}
\newcommand{\matr}[1]{\mspace{\matrspace}\underline{\mspace{-\matrspace}\smash[b]{\boldsymbol{#1}}\mspace{-\matrspace}}\mspace{\matrspace}}
\newcommand{\vect}[1]{{\boldsymbol{#1}}}
\newcommand{\dif}{\mathrm{d}}
\newcommand{\R}{\mathbb{R}}
\newcommand{\N}{\mathbb{N}}
\newcommand{\Z}{\mathbb{Z}}
\newcommand{\Q}{\mathbb{Q}}
\newcommand{\C}{\mathbb{C}}
\newcommand{\K}{\mathbb{K}}
\newcommand{\F}{\mathbb{F}}

% commands for probability
\newcommand{\Var}{\operatorname{Var}}
\newcommand{\Cov}{\operatorname{Cov}}
\newcommand{\Exp}{\operatorname{E}}
% \newcommand{\P}{\operatorname{P}} % this is already defined in amsmath/amsopn
\newcommand{\Prob}{\operatorname{P}}
\newcommand{\numof}{\ensuremath{\# \,}} % number of elements in a set
\newcommand{\blackheight}{\operatorname{bh}}

% \algnewcommand{\LeftComment}[1]{\(\triangleright\) #1}
\algnewcommand{\TO}{, \ldots ,}
\algnewcommand{\DOWNTO}{, \ldots ,}
\algnewcommand{\OR}{\vee}
\algnewcommand{\AND}{\wedge}
\algnewcommand{\NOT}{\neg}
\algnewcommand{\LEN}{\operatorname{len}}
\algnewcommand{\tru}{\ensuremath{\mathrm{\texttt{true}}}}
\algnewcommand{\fals}{\ensuremath{\mathrm{\texttt{false}}}}
\algnewcommand{\append}{\circ}
\algnewcommand{\nil}{\ensuremath{\mathrm{\textsc{nil}}}}
\algnewcommand{\red}{\ensuremath{\mathrm{\textsc{red}}}}
\algnewcommand{\black}{\ensuremath{\mathrm{\textsc{black}}}}
\algnewcommand{\gray}{\ensuremath{\mathrm{\textsc{gray}}}}
\algnewcommand{\white}{\ensuremath{\mathrm{\textsc{white}}}}
% \newcommand{\attribute}[1]{\ensuremath{\mathtt{#1}}}
\newcommand{\attrib}[2]{\ensuremath{#1\mathtt{.}\mathtt{#2}}} % object.field with field in math typewriter (like code)
\newcommand{\attribnormal}[2]{\ensuremath{#1\mathtt{.}#2}} 

\title{Probability and Measure}
\author{Fabian Bosshard}
\date{}
% \date{\today}

\usepackage[linktoc=none, pdfauthor={Fabian Bosshard}, pdftitle={USI - Probability and Measure - Course Notes}, pdfkeywords={USI, Probability and Measure, course notes, informatics}, colorlinks=false, pdfborder={0.0 0.0 0.0}, linkbordercolor={0 0.6 1}, urlbordercolor={0 0.6 1}, citebordercolor={0 0.6 1}]{hyperref}

\DeclareTOCStyleEntry[linefill=\dotfill]{tocline}{section}
\DeclareTOCStyleEntry[linefill=\dotfill]{tocline}{section}
\DeclareTOCStyleEntry[linefill=\dotfill]{tocline}{subsection}
\DeclareTOCStyleEntry[linefill=\dotfill]{tocline}{subsubsection}
\makeatletter
\newlength\FB@toclinkht
\newlength\FB@toclinkdp
\setlength\FB@toclinkht{.80\ht\strutbox}
\setlength\FB@toclinkdp{.40\dp\strutbox}

% --- ToC: add a target at the entry + one full-width clickable line ---
\let\FB@orig@contentsline\contentsline
\renewcommand*\contentsline[4]{%
  \begingroup
    \Hy@safe@activestrue
    \edef\Hy@tocdestname{#4}%
    \FB@orig@contentsline{#1}{%
      % target for "back to ToC"
      \Hy@raisedlink{\hyper@anchorstart{toc:\Hy@tocdestname}\hyper@anchorend}%
      \leavevmode
      \rlap{%
        \hyper@linkstart{link}{\Hy@tocdestname}%
          \raisebox{0pt}[\FB@toclinkht][\FB@toclinkdp]{%
            \hbox to \dimexpr\hsize-\parindent\relax{\hfil}%
          }%
        \hyper@linkend
      }%
      #2%
    }{#3}{#4}%
  \endgroup
}

% --- Headings: make NUMBER + TITLE one link back to the ToC entry ---
\let\FB@orig@sectionlinesformat\sectionlinesformat
\renewcommand{\sectionlinesformat}[4]{%
  \ifx\@currentHref\@empty
    \FB@orig@sectionlinesformat{#1}{#2}{#3}{#4}%
  \else
    \leavevmode
    \hyper@linkstart{link}{toc:\@currentHref}%
      % IMPORTANT: \hskip needs a DIMENSION only; #3 is text -> separate them
      \@hangfrom{\hskip #2\relax #3}{#4}%
    \hyper@linkend
    \par
  \fi
}
\makeatother

\usepackage[type={CC}, modifier={by}, version={4.0}]{doclicense}
% \usepackage{cleveref}

% after \usepackage{hyperref} and \numberwithin{equation}{section}
\makeatletter
% make hyperref anchors match that so they’re unique
\renewcommand{\theHequation}{\thesection.\arabic{equation}}
\renewcommand{\theHtheorem}{\thesection.\arabic{theorem}} % theorem is the base counter; all siblings share its counter
\renewcommand{\theHlemma}{\theHtheorem}
\renewcommand{\theHcorollary}{\theHtheorem}
\renewcommand{\theHproposition}{\theHtheorem}
\renewcommand{\theHconjecture}{\theHtheorem}
\renewcommand{\theHclaim}{\theHtheorem}
\renewcommand{\theHfact}{\theHtheorem}
\renewcommand{\theHdefinition}{\thesection.\arabic{definition}}
\renewcommand{\theHproblem}{\theHdefinition}
\renewcommand{\theHexercise}{\thesection.\arabic{exercise}} % Solution is numbered=no → no counter, no \theHsolution needed
\renewcommand{\theHremark}{\thesection.\arabic{remark}}
\renewcommand{\theHexample}{\theHremark}
\renewcommand{\theHcaution}{\theHremark}
\renewcommand{\theHsmallremark}{\theHremark}
\renewcommand{\theHdigression}{\theHremark}
\makeatother

% load glossaries *after* hyperref
\usepackage[acronym, nomain, toc, nonumberlist]{glossaries-extra}

% autoref names ~~~~~~~~~~~~~~~~~~~~~~~~~~~~~~~~
\renewcommand{\definitionautorefname}{Definition}
\renewcommand{\lemmaautorefname}{Lemma}
\renewcommand{\theoremautorefname}{Theorem}
\renewcommand{\corollaryautorefname}{Corollary}
\renewcommand{\propositionautorefname}{Proposition}
\renewcommand{\conjectureautorefname}{Conjecture}
\renewcommand{\claimautorefname}{Claim}
\renewcommand{\exampleautorefname}{Example}
\renewcommand{\remarkautorefname}{Remark}
\renewcommand{\cautionautorefname}{Caution}
\renewcommand{\digressionautorefname}{Digression}
\makeatletter
\def\thmt@dummyctrautorefname~#1\null{Proof\null}
\makeatother
\makeatletter
\patchcmd{\ALG@step}{\addtocounter{ALG@line}{1}}{\refstepcounter{ALG@line}}{}{}
\newcommand{\ALG@lineautorefname}{Line}
\makeatother
\newcommand{\algorithmautorefname}{Algorithm}
\AtBeginDocument{\renewcommand{\sectionautorefname}{Section}\renewcommand{\subsectionautorefname}{Section}\renewcommand{\subsubsectionautorefname}{Section}}

% \setcounter{tocdepth}{1}

\begin{document}
\pagenumbering{roman}
\maketitle
\tableofcontents
% \section*{Preface}
% \addcontentsline{toc}{section}{Preface}

% \doclicenseThis

\let\emphoriginal\emph           % save original meaning
\renewcommand{\emph}[1]{\textcolor{black}{#1}}
\printbibliography[heading=bibintoc,title={References}]
\let\emph\emphoriginal           % restore original meaning



\clearpage
\pagenumbering{arabic}



\section{Basic Stuff}
Contingency table: 
displays the multivariate frequency distribution of the variables





\subsection{Conditional Probability and Bayes' Theorem}

\subsubsection*{Events}

\paragraph{Conditional probability}

\[
\Prob(A \mid B)
=
\frac{\Prob(A \cap B)}{\Prob(B)}
\]


\paragraph{Bayes' theorem}

\[
\Prob(A \mid B)
=
\frac{\Prob(B \mid A)\Prob(A)}
{\Prob(B)}
\]



\paragraph{Total probability}

Let $(A_i)_i$ be a partition of $\Omega$. Then

\[
\begin{aligned}
\Prob(B)
&=
\sum_i \Prob(B \cap A_i)
\\
&=
\sum_i \Prob(B \mid A_i)\Prob(A_i)
\end{aligned}
\]


\bigskip

\subsubsection*{From Events to Random Variables}

If $X,Y$ are random variables, the same structure appears at the level of distributions.

\bigskip

\subsubsection*{Random Variables}

\[
\begin{array}{@{}p{0.46\textwidth} @{\hspace{0.06\textwidth}} p{0.46\textwidth}@{}}
\textbf{Discrete case (pmf)} 
&
\textbf{Continuous case (pdf)}
\\[1em]

\textbf{Joint distribution}
&
\textbf{Joint distribution}
\\[0.5em]

\[
p_{X,Y}(x,y)
=
\Prob(X=x, Y=y)
\]
&
\[
f_{X,Y}(x,y)
\]

\\[2em]

\textbf{Marginal distribution}
&
\textbf{Marginal distribution}
\\[0.5em]

\[
\begin{aligned}
p_Y(y)
&=
\sum_x p_{X,Y}(x,y)
\end{aligned}
\]
&
\[
\begin{aligned}
f_Y(y)
&=
\int f_{X,Y}(x,y)\,\dif x
\end{aligned}
\]

\\[2em]

\textbf{Conditional distribution}
&
\textbf{Conditional distribution}
\\[0.5em]

\[
p_{X\mid Y}(x \mid y)
=
\frac{p_{X,Y}(x,y)}{p_Y(y)}
\]
&
\[
f_{X\mid Y}(x \mid y)
=
\frac{f_{X,Y}(x,y)}{f_Y(y)}
\]

\\[2em]

\textbf{Bayes' theorem}
&
\textbf{Bayes' theorem}
\\[0.5em]

\[
\begin{aligned}
p_{X\mid Y}(x \mid y)
&=
\frac{p_{Y\mid X}(y \mid x)p_X(x)}
{\sum_u p_{Y\mid X}(y \mid u)p_X(u)}
\end{aligned}
\]
&
\[
\begin{aligned}
f_{X\mid Y}(x \mid y)
&=
\frac{f_{Y\mid X}(y \mid x)f_X(x)}
{\int f_{Y\mid X}(y \mid u)f_X(u)\,\dif u}
\end{aligned}
\]

\end{array}
\]


\section{Probability Spaces}

Let \(\Omega\) be a \hl{non-empty} set.
We call \(\Omega\) the sample space, which specifies the possible outcomes of a random trial.

Based on the \emph{elementary outcomes} \(\omega \in \Omega\), one can also consider more complex \emph{events} \(A \subseteq \Omega\), which are based on the usual set operations.

If \(P\) is a property of elements of \(\Omega\), then the set
\[
\{\omega \in \Omega: P(\omega)\} \subset \Omega
\]
is the subset of all elements \(\omega \in \Omega\) for which \(P\) is true. 
Furthermore, for \(A \subset \Omega\), we write
\[
A^{\complement} := \{\omega \in \Omega: \omega \notin A\}
\]
for the \emph{complement} of \(A\) in \(\Omega\). 
The complement of \(\Omega\) is the empty set \(\emptyset\).
Given two sets \(A, B \subset \Omega\), we further define their \emph{union} by
\[
A \cup B:=\{\omega \in \Omega: \omega \in A \text { or } \omega \in B\}
\]
their \emph{intersection} by
\[
A \cap B:=\{\omega \in \Omega: \omega \in A \text { and } \omega \in B\}
\]
and their \emph{difference} by
\[
A \backslash B:=A \cap B^{\complement} 
\]


The collection of all possible subsets of \(\Omega\) is called the \emph{power set} of \(\Omega\) and is defined as
\[
2^{\Omega}:=\{A: A \subseteq \Omega\}
\]




We are often interested in certain subsets of the sample space. 
To work with these subsets, we define a collection of events that is closed under the usual set operations.
\begin{definition}[Algebra]\label{def:algebra}
A set \(\mathcal{F} \subseteq 2^{\Omega}\) is called an \emph{algebra} on \(\Omega\) if
\begin{enumerate}[label=(\roman*)]
\item \(\Omega \in \mathcal{F}\) \label{def:algebra-contains-samplespace}
\item \(A \in \mathcal{F} \Rightarrow A^{\complement} \in \mathcal{F}\) \label{def:algebra-closed-under-complement}
\item \(A, B \in \mathcal{F} \Rightarrow A \cup B \in \mathcal{F}\) \label{def:algebra-closed-under-union}
\end{enumerate}
The sets \(A \in \mathcal{F}\) are called \emph{events}.
\end{definition}

The smallest algebra is \(\{\emptyset, \Omega\}\).
Since for any algebra \(\mathcal{F} \subseteq 2^{\Omega}\), \(2^\Omega\) is the largest algebra.
Note that any algebra is also closed under intersection, since \(A \cap B = (A^{\complement} \cup B^{\complement})^{\complement}\) (De Morgan).





\begin{example}[Different information $\Rightarrow$ different algebras]\label{ex:info-algebras}
Alice generates two random bits. The sample space therefore is
\(
\Omega=\{00,01,10,11\}
\).

Alice observes the exact outcome, hence her observable events are all subsets of $\Omega$: 
\[
\mathcal F_A = 2^\Omega
\]

Bob is only told the \emph{number of ones} in the two bits. This induces the partition
\[
\Pi_B=\bigl\{\{00\},\{01,10\},\{11\}\bigr\}
\]
of the sample space \(\Omega\).
Therefore Bob can only observe events that are unions of blocks of $\Pi_B$, i.e.\ the algebra generated by $\Pi_B$:
\[
\mathcal F_B
=
\Bigl\{
\emptyset,\{00\},\{01,10\},\{11\},\{00,11\},\{00,01,10\},\{01,10,11\},\Omega
\Bigr\}
\]

Cindy is only told whether the bits are \emph{equal or different}. This induces the coarser partition
\[
\Pi_C=\bigl\{\{00,11\},\{01,10\}\bigr\},
\]
hence her observable events form the algebra:
\[
\mathcal F_C=\bigl\{\emptyset,\{00,11\},\{01,10\},\Omega\bigr\}
\]

We have
\[
\mathcal F_A \supset \mathcal F_B \supset \mathcal F_C
\]
expressing that more information corresponds to a finer partition and thus a larger algebra of observable events.
\end{example}



Often, we encounter situations with infinitely many events and we want to assign probabilities to the case that all of them occur simultaneously.
This is possible, as long as the number of events involved is at most countably infinite.

\begin{definition}[\(\sigma\)-algebra]\label{def:sigma-algebra}
A set \(\mathcal{F} \subseteq 2^{\Omega}\) is called a \emph{\(\sigma\)-algebra} on \(\Omega\) if
\begin{enumerate}[label=(\roman*)]
\item \(\mathcal{F}\) is an algebra on \(\Omega\) \label{def:sigma-algebra-is-algebra}
\item \(A_1, A_2, \ldots \in \mathcal{F} \implies \bigcup_{i=1}^\infty A_i \in \mathcal{F}\) \label{def:sigma-algebra-closed-under-countable-union}
\end{enumerate}
The sets in \(\mathcal{F}\) are called \emph{measurable} and the tuple \((\Omega, \mathcal{F})\) is called a \emph{measurable space}.
\end{definition}






\end{document}